\section{Conformance and CI}
\label{sec:conformance}

The cross-language conformance harness is the single quality signal that catches polyglot drift. This section documents how it works at the umbrella level; per-project KAT vector inventories live in the per-project whitepapers.

\subsection{Repository layout}

\begin{center}
\small
\begin{tabular}{@{}ll@{}}
\toprule
\textbf{Path} & \textbf{Contents} \\
\midrule
\texttt{tessera-id/conformance/vectors/}     & TesseraID KAT files (decode, encode/decode round-trip, invalid, HLC overflow, min timestamp) \\
\texttt{tessera-hash/conformance/vectors/}   & TesseraHash KAT files (decode, hash\_only, encode/decode B64URL, encode/decode Crockford, verify, invalid) \\
\texttt{tessera-seal/conformance/vectors/}  & TesseraSeal KAT files (decrypt, ad\_binding, commitment, kyberslash\_regression, invalid, peek\_header, hybrid\_*) \\
\texttt{tools/conformance/run.sh}      & Single cross-language entry point \\
\texttt{tools/conformance/<lang>/}     & Per-language runner sources and binaries \\
\texttt{.github/workflows/conformance.yml} & Branch-protection gate on \texttt{main} \\
\texttt{.github/workflows/tessera\_seal.yml}     & TesseraSeal-only matrix CI (TSEAL-00023) \\
\bottomrule
\end{tabular}
\end{center}

\subsection{The hybrid suite \texorpdfstring{\code{0x06}}{0x06} skip discipline}

\tesseraseal{}'s hybrid suite \code{0x06} (TSEAL-00021, v1.2) is a polyglot KAT contract that some language ports have not yet caught up to as of v1.3. The harness filters \texttt{hybrid\_*.json} vector files per language via the \texttt{hybrid\_supported()} function in \texttt{tools/conformance/run.sh}:

\begin{lstlisting}[language=bash]
hybrid_supported() {
    case "$1" in
        python) return 0 ;;   # TSEAL-00021, v1.2
        rust)   return 0 ;;   # TSEAL-00028, v1.3
        go)     return 0 ;;   # TSEAL-00029, v1.3
        *)      return 1 ;;   # C / C++ / TS / Java / CPython ext: v1.4 backlog
    esac
}
\end{lstlisting}

Languages without hybrid skip the \texttt{hybrid\_*.json} files entirely and MUST reject suite \code{0x06} as \code{UNKNOWN\_SUITE} for any non-hybrid envelope they encounter. When porting hybrid to a new language: (i) implement hybrid in that language, (ii) update the language's conformance runner to accept 96-byte hybrid seeds (the standard runners hard-coded \texttt{len(seed) != 64}), (iii) flip the language's \texttt{hybrid\_supported()} entry, (iv) verify the language's PASS count rises by 9 (one per record across the three hybrid KAT files).

\subsection{Latest cross-language state (2026-05-21, Apple M4 Max)}

\begin{center}
\small
\begin{tabular}{@{}lrrl@{}}
\toprule
\textbf{Language} & \textbf{PASS} & \textbf{FAIL} & \textbf{Notes} \\
\midrule
Python   & 265 & 0 & TesseraID + TesseraHash + TesseraSeal v1.x including hybrid \code{0x06}. \\
C        & 256 & 0 & TesseraSeal hybrid skipped; TSEAL-00030 backlog. \\
Go       & 265 & 0 & TesseraSeal hybrid via stdlib \texttt{crypto/ecdh} (TSEAL-00029). \\
Rust     & 265 & 0 & TesseraSeal hybrid via \texttt{x25519-dalek} 2 (TSEAL-00028). \\
\bottomrule
\end{tabular}
\end{center}

C++, TypeScript, Java, and CPython ext consume TesseraSeal vectors via the Python runner fallback (\texttt{tessera\_seal\_runner\_cmd} empty case in \texttt{tools/conformance/run.sh}) and skip \texttt{hybrid\_*.json}. Their per-language test suites still cover TesseraID and TesseraHash directly.

\subsection{How CI gates the work}

Every PR to \texttt{main} runs the harness across all available language toolchains in the GitHub Actions matrix. The harness is non-flaky by design: KAT vectors are deterministic, the per-language runners do not call out to the network, and the toolchains are pinned. A red harness means one of two things:

\begin{itemize}
\item A language port produced a different byte sequence for the same input. The port is wrong; fix the port.
\item The KAT vector itself changed. Re-running \texttt{<project>/conformance/scripts/generate.py} produces non-identical JSON. The Python reference's behavior changed; this is a spec change, requires a release-notes entry, and all other language ports must be updated to match before the PR merges.
\end{itemize}

The harness is what makes the spec-first methodology load-bearing. Without it, ``Python is right'' is just a slogan; with it, a language port that disagrees with Python is detectable in seconds.
